\documentclass[exercice]{thedocument}%
\startdatakeys
\source{originale}
\BO{2026}
\cycle{4}
\competencesDuSocle{CA3,RA3,CO1}
\objectifsApprentissage{C4A1ea,C4A1eg}
\motsCles{calcul littéral, successeur, prédécesseur, nombre pair}
\enddatakeys
\begin{document}%
Soit $n$ un nombre entier pair. Traduire les phrases suivantes en notations mathématiques~:
\begin{enumerate}
    \item \begin{enumerate}
        \item Le successeur de $n$ (le nombre entier juste après $n$).
        \item Le prédecesseur de $n$ (le nombre entier juste avant $n$).
        \item La moitié de $n$.
        \item Le double de $n$.
        \item Le prédecesseur de la moitié de $n$.
        \item Le successeur du successeur de $n$.
    \end{enumerate}
    \item Démontrer par un calcul que le double du prédecesseur de la moitié d'un nombre
    est égal au prédécesseur du prédécesseur de ce nombre.
\end{enumerate}
\end{document}
